% questions.tex —— 2006 年·北京卷·文科 · 题目内容片段（无答案）
% 由 tools/generate_projects.py 依来源总账自 2006/2006_数学真题_北京文_试题.md 生成，勿手改（改源 MD 后重跑）


\begin{gktip}
\textbf{来源边界}：题签依据固定 PDF 首页照录。题面按固定 PDF 第1—2页转录，印刷页码为282—283。该固定副本的发布者、出版者、版次、来源链及官方身份均未核实；\texttt{ready} 仅为来源清单中的文件状态，不代表来源认证。固定 PDF 未见答案、解析或评分信息。
\end{gktip}

\gkpart{一、选择题}

\begin{gktip}
以下第1页（印刷页282）含第1—8题。
\end{gktip}

\begin{question}
\gkti{1}设集合 $A=\{x\mid 2x+1<3\}$，$B=\{x\mid -3<x<2\}$，则 $A\cap B$ 等于\mbox{（　　）}。

（A）$\{x\mid -3<x<1\}$　　（B）$\{x\mid 1<x<2\}$

（C）$\{x\mid x>-3\}$　　（D）$\{x\mid x<1\}$

\end{question}
\begin{question}
\gkti{2}函数 $y=1+\cos x$ 的图象\mbox{（　　）}。

（A）关于 $x$ 轴对称　　（B）关于 $y$ 轴对称

（C）关于原点对称　　（D）关于直线 $x=\dfrac{\pi}{2}$ 对称

\end{question}
\begin{question}
\gkti{3}若 $\vec a$ 与 $\vec b-\vec c$ 都是非零向量，则“$\vec a\cdot\vec b=\vec a\cdot\vec c$”是“$\vec a\perp(\vec b-\vec c)$”的\mbox{（　　）}。

（A）充分而不必要条件　　（B）必要而不充分条件

（C）充分必要条件　　（D）既不充分也不必要条件

\end{question}
\begin{question}
\gkti{4}在 $1,2,3,4,5$ 这五个数字组成的没有重复数字的三位数中，各位数字之和为奇数的共有\mbox{（　　）}。

（A）36个　　（B）24个　　（C）18个　　（D）6个

\end{question}
\begin{question}
\gkti{5}已知
\nopagebreak[4]
\[
f(x)=\begin{cases}
(3-a)x-4a, & x<1,\\
\log_a x, & x\ge 1,
\end{cases}
\]
是 $(-\infty,+\infty)$ 上的增函数，那么 $a$ 的取值范围是\mbox{（　　）}。

（A）$(1,+\infty)$　　（B）$(-\infty,3)$　　（C）$\left[\dfrac35,3\right)$　　（D）$(1,3)$

\end{question}
\begin{question}
\gkti{6}如果 $-1,a,b,c,-9$ 成等比数列，那么\mbox{（　　）}。

（A）$b=3,\ ac=9$　　（B）$b=-3,\ ac=9$

（C）$b=3,\ ac=-9$　　（D）$b=-3,\ ac=-9$

\end{question}
\begin{question}
\gkti{7}设 $A,B,C,D$ 是空间四个不同的点，在下列命题中，不正确的是\mbox{（　　）}。

（A）若 $AC$ 与 $BD$ 共面，则 $AD$ 与 $BC$ 共面。

（B）若 $AC$ 与 $BD$ 是异面直线，则 $AD$ 与 $BC$ 是异面直线。

（C）若 $AB=AC$，$DB=DC$，则 $AD=BC$。

（D）若 $AB=AC$，$DB=DC$，则 $AD\perp BC$。

\end{question}
\begin{question}
\gkti{8}如图为某三岔路口交通环岛的简化模型。在某高峰时段，单位时间进出路口 $A,B,C,$ 的机动车辆数如图所示，图中 $x_1,x_2,x_3$ 分别表示该时段单位时间通过路段 $\widehat{AB},\widehat{BC},\widehat{CA}$ 的机动车辆数（假设：单位时间内，在上述路段中，同一路段上驶入与驶出的车辆数相等），则 $x_1,x_2,x_3$ 的大小关系是\mbox{（　　）}。

（A）$x_1>x_2>x_3$　　（B）$x_1>x_3>x_2$

（C）$x_2>x_3>x_1$　　（D）$x_3>x_2>x_1$

\begin{gktip}
【原卷配图位置：固定 PDF 第1页，印刷页282，Q8。图含路口 $A,B,C$、路段流量 $x_1,x_2,x_3$、进出车辆数标注及方向箭头；题面图稿待重绘。】
\end{gktip}

\end{question}
\gkpart{二、填空题}

\begin{gktip}
第9—14题位于固定 PDF 第1页（印刷页282）；同页的第15—17题属于下节解答题。
\end{gktip}

\begin{question}
\gkti{9}若三点 $A(2,2)$、$B(a,0)$、$C(0,4)$ 共线，则 $a$ 的值等于\gkblank{}。

\end{question}
\begin{question}
\gkti{10}在 $\left(x-\dfrac2x\right)^7$ 的展开式中，$x^3$ 的系数是\gkblank{}。（用数字作答）

\end{question}
\begin{question}
\gkti{11}已知函数 $f(x)=a^x-4a+3$ 的反函数的图象经过点 $(-1,2)$，那么 $a$ 的值等于\gkblank{}。

\end{question}
\begin{question}
\gkti{12}已知向量 $\vec a=(\cos\alpha,\sin\alpha)$，$\vec b=(\cos\beta,\sin\beta)$，且 $\vec a\ne\pm\vec b$，那么 $\vec a+\vec b$ 与 $\vec a-\vec b$ 的夹角的大小是\gkblank{}。

\end{question}
\begin{question}
\gkti{13}在 $\triangle ABC$ 中，$\angle A,\angle B,\angle C$ 所对的边长分别为 $a,b,c$。若 $\sin A:\sin B:\sin C=5:7:8$，则 $a:b:c=$\gkblank{}，$\angle B$ 的大小是\gkblank{}。

\end{question}
\begin{question}
\gkti{14}已知点 $P(x,y)$ 的坐标满足条件
\nopagebreak[4]
\[
\begin{cases}
x+y\le 4,\\
y\ge x,\\
x\ge 1,
\end{cases}
\]
点 $O$ 为坐标原点，那么 $|PO|$ 的最小值等于\gkblank{}，最大值等于\gkblank{}。

\end{question}
\gkpart{三、解答题}

\begin{question}
\gkti{15}已知函数
\nopagebreak[4]
\[
f(x)=\frac{1-\sin 2x}{\cos x}.
\]

（1） 求 $f(x)$ 的定义域；

（2） 设 $\alpha$ 是第四象限的角，且 $\tan\alpha=-\dfrac43$，求 $f(\alpha)$ 的值。

\end{question}
\begin{question}
\gkti{16}已知函数 $f(x)=ax^3+bx^2+cx$ 在点 $x_0$ 处取得极大值5，其导函数 $y=f'(x)$ 的图象经过点 $(1,0),(2,0)$，如图所示，求：

（1） $x_0$ 的值；

（2） $a,b,c$ 的值。

\begin{gktip}
【原卷配图位置：固定 PDF 第1页，印刷页282，Q16。图为导函数 $y=f'(x)$ 的图象，标有 $x=1,2$ 处的横轴交点；题面图稿待重绘。】
\end{gktip}

\end{question}
\begin{question}
\gkti{17}如图，$ABCD-A_1B_1C_1D_1$ 是正四棱柱。

（1） 求证：$BD\perp$ 平面 $ACC_1A_1$；

（2） 若二面角 $C_1-BD-C$ 的大小为 $60^\circ$，求异面直线 $BC_1$ 与 $AC$ 所成角的大小。

\begin{gktip}
【原卷配图位置：固定 PDF 第1页，印刷页282，Q17；题面空间图待重绘。】
以下第18—20题位于固定 PDF 第2页（印刷页283）。
\end{gktip}

\end{question}
\begin{question}
\gkti{18}某公司招聘员工，指定三门考试课程，有两种考试方案。

方案一：考试三门课程，至少有两门及格为考试通过；

方案二：在三门课程中随机选取两门，这两门都及格为考试通过。

假设某应聘者对三门指定课程考试及格的概率分别是 $0.5,0.6,0.9$，且三门课程考试是否及格相互之间没有影响。

（1） 该应聘者用方案一考试通过的概率；

（2） 该应聘者用方案二考试通过的概率。

\end{question}
\begin{question}
\gkti{19}椭圆 $C:\dfrac{x^2}{a^2}+\dfrac{y^2}{b^2}=1\ (a>b>0)$ 的两个焦点为 $F_1,F_2$，点 $P$ 在椭圆 $C$ 上，且 $PF_1\perp F_1F_2$，$|PF_1|=\dfrac43$，$|PF_2|=\dfrac{14}3$。

（1） 求椭圆 $C$ 的方程；

（2） 若直线 $l$ 过圆 $x^2+y^2+4x-2y=0$ 的圆心 $M$，交椭圆 $C$ 于 $A,B$ 两点，且 $A,B$ 关于点 $M$ 对称，求直线 $l$ 的方程。

\end{question}
\begin{question}
\gkti{20}设等差数列 $\{a_n\}$ 的首项 $a_1$ 及公差 $d$ 都为整数，前 $n$ 项和为 $S_n$。

（1） 若 $a_{11}=0$，$S_{14}=98$，求数列 $\{a_n\}$ 的通项公式；

（2） 若 $a_1\ge 6$，$a_{11}>0$，$S_{14}\ge 77$，求所有可能的数列 $\{a_n\}$ 的通项公式。

\par\medskip

\end{question}